

How does delegating grading authority to schools change the allocation of students across universities? I study the cancellation of standardized A-level exams in the United Kingdom in 2020--2021, when teacher-assessed grades replaced centralized exams after universities had already made conditional offers. I identify each school's grading discretion by measuring deviations in average grades from the school's historical levels, which were mechanically fixed by central graders in normal years. I then estimate how grade inflation shifted university composition using the average inflation of each program's applicant pool as an instrument. I find that independent schools adopted looser grading standards than publicly funded schools, increasing their students' representation at elite programs; the gap widened by roughly 50\% in the second year as schools further loosened standards while universities tightened entry requirements. Grade inflation was nearly twice as large in non-STEM subjects, and within courses, schools assigned higher grades to female, white, and high-SES students---with larger gaps at independent schools. By tightening admissions to limit enrollment, universities replaced high-achieving, socioeconomically disadvantaged students with lower-achieving students from wealthier backgrounds, displacing disadvantaged students to lower-ranked institutions. The findings show that the interaction between school grading discretion and university admission responses reshuffled access to selective universities in a regressive direction, rather than diversifying it.